% !TEX root = ../libretto.tex
% Sezione: Canti finali
\momento{Canti finali}
\begin{songs}{titoli}

% Fonte: LibrettoAccordi.pdf (Canti Coro), J26.
% Nell'elenco: Finale, Ordinario.
\beginsong{Al centro del presente}[by={Gen Rosso}]
\beginchorus
Oh, oh, oh, vivi il momento che va
cogli l'attimo adesso, eh, eh eh
Oh, oh, oh, vivi il momento che va
cogli l'attimo adesso.
\endchorus
\beginverse
Lui va a scuola, ma non sa perché, butta i giorni nella noia;
vive dentro una fotografia inseguendo un sogno fragile.
Lui non ha tempo per nessuno, sta viaggiando nel futuro.
\endverse
\beginverse
Lei ripensa ai passi che non ha sta vagando nel passato;
ha negli occhi un che di arido, se potesse cambierebbe tutto di sé.
Ma ormai l'acqua che è passata non fa più girar la ruota.
\endverse
\beginverse
Al centro di questo presente scorre l'eterno nelle cose.
E se in quest'attimo tu ami incontri Dio nella tua vita.
\endverse
\endsong

% Fonte: LibrettoAccordi.pdf (Canti Coro), Q1.
% Nell'elenco: Finale, Ordinario.
\beginsong{Alma misionera}[by={P. E. García Vélez}]
\beginverse
Señor, toma mi vida nueva,
antes de que la espera desgaste años en mí.
Estoy dispuesto a lo que quieras,
no importa lo que sea, Tú llámame a servir.
\endverse
\beginchorus
Llévame donde los hombres necesiten Tus palabras,
necesiten mis ganas de vivir.
Donde falte la esperanza, donde falte la alegría
simplemente por no saber de ti.
\endchorus
\beginverse
Te doy mi corazón sincero,
para gritar sin miedo lo hermoso que es tu amor.
Señor, tengo alma misionera,
condúceme a la tierra que tenga sed de ti.
\endverse
\beginverse
Así, en marcha iré cantando,
por calles predicando tu grandeza, Señor.
Tendré mis manos sin cansancio,
tu historia entre mis labios y fuerza en la oración.
\endverse
\endsong

% Fonte: LibrettoAccordi.pdf (Canti Coro), K5.
% Nell'elenco: senza tipologia (nel libretto accordi tra i canti finali).
\beginsong{E la strada si apre}[by={Gen Arcobaleno}]
\beginverse
Raggio che buca le nubi ed è già cielo aperto,
acqua che scende decisa scavando da sé l'argine per la vita,
la traiettoria di un volo che sull'orizzonte di sera,
tutto di questa natura è una strada per sé.
\endverse
\beginverse
Attimo che segue attimo, un salto nel tempo,
passo di un mondo che tende oramai all'unità, che non è più domani,
usiamo allora queste mani, scaviamo a fondo nel cuore,
solo scegliendo l'amore il mondo vedrà.
\endverse
\beginchorus
Che la strada si apre, passo dopo passo, ora su questa strada noi.
E si spalanca un cielo, un mondo che rinasce, si può vivere nell'unità.
\endchorus
\beginverse
Nave che segue una rotta in mezzo alle onde,
uomo che s'apre la strada in una giungla di idee, seguendo sempre il sole,
quando si sente assetato deve raggiungere l'acqua,
sabbia che nella risacca ritorna al mare.
usiamo allora queste mani, scaviamo a fondo nel cuore,
solo scegliendo l'amore il mondo vedrà.
\endverse
\endsong

% Fonte: LibrettoAccordi.pdf (Canti Coro), K10.
% Nell'elenco: Inno GMG 2000.
\beginsong{Emmanuel}[by={M. Mammoli}]
\beginverse
Dall'orizzonte una grande luce viaggia nella storia
e lungo gli anni ha vinto il buio facendosi memoria,
e illuminando la nostra vita chiaro ci rivela
che non si vive se non si cerca la Verità.
\endverse
\beginverse
Da mille strade arriviamo a Roma sui passi della fede,
sentiamo l'eco della Parola che risuona ancora
da queste mura, da questo cielo per il mondo intero:
è vivo oggi, è l'Uomo Vero, Cristo tra noi.
\endverse
\beginchorus
Siamo qui, sotto la stessa luce, sotto la sua croce, cantando ad una voce:
è l'Emmanuel, Emmanuel, Emmanuel.
È l'Emmanuel, Emmanuel.
\endchorus
\beginverse
Dalla città di chi ha versato il sangue per amore,
ed ha cambiato il vecchio mondo vogliamo ripartire.
Seguendo Cristo, insieme a Pietro, rinasce in noi la fede,
Parola viva che ci rinnova e cresce in noi.
\endverse
\beginverse
Un grande dono che Dio ci ha fatto è Cristo il suo Figlio,
e l'umanità è rinnovata, è in Lui salvata.
È vero uomo, è vero Dio, è il Pane della Vita,
che ad ogni uomo, ai suoi fratelli ridonerà.
\endverse
\beginverse
La morte è uccisa, la vita ha vinto, è Pasqua in tutto il mondo,
un vento soffia in ogni uomo lo Spirito fecondo,
Che porta avanti nella storia la Chiesa sua sposa,
sotto lo sguardo di Maria Comunità.
\endverse
\endsong

% Fonte: LibrettoAccordi.pdf (Canti Coro), J32.
% Nell'elenco: Finale, Ordinario.
\beginsong{Il mio spirito in festa}[by={M. Balduzzi (CD Xaipe)}]
\beginverse
Nessuno mai ha guardato a me,
con gli occhi nuovi del mattino,
nessuno mai ha pensato a me,
come ha pensato il mio Signore,
ha pensato il mio Signore.
\endverse
\beginchorus
L'anima mia esulta,
e il mio spirito in festa canta per te.
\endchorus
\beginverse
Nessuno mai solo resterà,
in ogni angolo del mondo,
nessuno mai ci separerà,
dalla presenza del tuo amore,
la presenza del tuo amore.
\endverse
\endsong

% Fonte: LibrettoAccordi.pdf (Canti Coro), J11.
% Nell'elenco: Finale, Ordinario.
\beginsong{Inno del missionario}
\beginverse
Signor, prendi la mia vita,
prima che l'attesa consumi anni in me.
Io sono aperto a ciò che mi chiedi,
qualunque cosa sia, Tu chiamami a servir.
\endverse
\beginchorus
Portami là dove all'uomo manchino le Tue parole
e la voglia di vivere con Te.
Dove manchi la speranza, dove manchi l'allegria
solamente perché non san di Te.
\endchorus
\beginverse
Ti do il mio cuor sincero,
per dire con coraggio quanto è bello il tuo amor.
Signor, ho un cuore missionario,
conducimi alla terra che si disseti in Te.
\endverse
\beginverse
Così io partirò cantando,
a tutti predicando la Tua immensità.
Sarò Tuo servo senza sosta,
il Tuo predicatore, la forza del Tuo amor.
\endverse
\endsong

% Fonte: LibrettoAccordi.pdf (Canti Coro), K7.
% Nell'elenco: Finale, Ordinario.
\beginsong{Le tue meraviglie}[by={M. Balduzzi, C. Casucci (CD Verbum Panis)}]
\beginchorus
Ora lascia o Signore che io vada in pace perché ho visto le tue meraviglie.
Il tuo popolo in festa per le strade correrà a portare le tue meraviglie.
\endchorus
\beginverse
La tua presenza ha riempito d'amore le nostre vite e le nostre giornate.
In te una sola anima, un solo cuore siamo noi
con te la luce risplende, splende più chiara che mai.
\endverse
\beginverse
La tua presenza ha inondato d'amore le nostre vite e le nostre giornate.
Fra la tua gente resterai, per sempre vivo in mezzo a noi
fino ai confini del tempo, così ci accompagnerai.
\endverse
\endsong

% Fonte: LibrettoAccordi.pdf (Canti Coro), A12. Le strofe "Ti ringraziamo..." sono parlate:
% tenute nel libretto su indicazione dell'utente.
% Nell'elenco: Finale, Ordinario.
\beginsong{Lodate Dio}[by={Gen Rosso}]
\beginchorus
Lodate Dio, cieli immensi ed infiniti.
Lodate Dio, cori eterni d'angeli.
Lodate Dio, santi del suo Regno.
Lodatelo uomini, Dio vi ama.
Lodatelo uomini, Dio è con voi.
\endchorus
\beginverse
Ti ringraziamo, Dio nostro Padre
perché sei Amore. Tu ci hai fatto dono della vita
e ci hai creati per essere figli tuoi.
Ti ringraziamo perché ci fai partecipi
della tua opera creatrice
dandoci un mondo da plasmare
con le nostre mani.
\endverse
\beginverse
Ti ringraziamo, Signore Gesù,
perché per amore nostro sei venuto in questo mondo
per riscattarci da ogni male e riportarci al Padre.
Ti ringraziamo perché sei rimasto
in mezzo a noi per sempre
per far di tutti gli uomini una sola famiglia,
un corpo solo in te.
\endverse
\beginverse
Ti ringraziamo, Spirito d'amore,
perché rinnovi la faccia della terra.
Tu dai luce e conforto ad ogni cuore.
Ti ringraziamo perché con i tuoi santi doni
ci dai la forza di avanzare nel nostro cammino
per giungere uniti alla gioia della tua casa.
\endverse
\endsong

% Fonte: LibrettoAccordi.pdf (Canti Coro), A14.
% Nell'elenco: Finale, Ordinario.
\beginsong{Musica di festa}[by={E. Sands, P. Jotti}]
\beginverse
Cantate al Signore un cantico nuovo, grande è la sua gloria!
Grande la sua forza, grande la sua pace, grande la sua santità!
\endverse
\beginchorus
In tutta la terra, popoli del mondo, gridate la sua fedeltà!
Musica di festa, musica di lode, musica di libertà!
\endchorus
\beginverse
Agli occhi del mondo ha manifestato la sua salvezza!
Per questo si canti, per questo si danzi, per questo si celebri!
\endverse
\beginverse
Con l'arpa ed il corno, con timpani e flauti, con tutta la voce,
canti di dolcezza, canti di salvezza, canti d'immortalità.
\endverse
\beginverse
I fiumi ed i monti battono le mani davanti al Signore.
La sua giustizia giudica la terra, giudica le genti!
\endverse
\beginverse
Al Dio che ci salva, gloria in eterno. Amen! Alleluja!
Gloria a Dio Padre, gloria a Dio Figlio, gloria a Dio Spirito!
\endverse
\endsong

% Fonte: testo fornito dall'utente.
% Nell'elenco: Comunione; l'utente lo indica come Finale.
\beginsong{Nelle mani tue}[by={Gen Verde}]
\beginverse
Sopra un mare di luce,
il vento nelle nostre vele,
sulla rotta di un sogno
nell'immenso blu.
Navighiamo sicuri,
siamo uniti, siamo insieme
e il sole splende sempre più.
\endverse
\beginchorus
E tutto parla di te,
del tuo amore per noi,
e noi crediamo in te, Signore,
sei tu il nostro Dio:
noi siamo nelle mani tue.
\endchorus
\beginverse*
We are, we are, we are,
we are in your hands. \rep{2}
\endverse
\beginverse
Ora il vento è più forte,
si sollevano le onde,
è calato il sole,
c'è solo freddo e buio.
Dove sei, o Signore?
Siamo persi nella notte,
ci puoi salvare solo tu.
\endverse
\beginchorus
E noi guardiamo a te,
al tuo amore per noi,
e noi crediamo in te, Signore,
sei tu il nostro Dio:
noi siamo nelle mani tue.
\endchorus
\beginverse
E quando il mare infuria
e ci prende la paura
e intorno a noi c'è solo oscurità,
tu cammini sulle acque
e ci chiami verso te.
Signore, sei davvero tu? \rep{2}
\endverse
\beginchorus
E tutto parla di te,
del tuo amore per noi,
cammineremo sulle acque
perché tu sei con noi:
noi siamo nelle mani tue.
E noi guardiamo a te,
al tuo amore per noi,
cammineremo sulle acque
perché tu sei con noi:
noi siamo nelle mani tue.
\endchorus
\beginverse*
We are, we are, we are,
we are in your hands. \rep{2}
\endverse
\beginverse*
Cammineremo sulle acque
perché tu sei con noi:
noi siamo nelle mani tue. \rep{2}
\endverse
\endsong

% Fonte: LibrettoAccordi.pdf (Canti Coro), J1.
% Nell'elenco: Finale, Ordinario.
\beginsong{Nozze a Cana}[by={RnS}]
\beginchorus
Danzano con gioia le figlie di Israele,
le nozze sono pronte per il figlio del re.
Alleluja, alleluja per il figlio del re. \rep{2}
\endchorus
\beginverse
A Cana in Galilea nel nome del Signore
gli sposi han giurato amore e fedeltà.
Coi cembali e coi flauti le danze fan corona,
il vino sulla mensa il cuor rallegrerà.
\endverse
\beginverse
Nel mezzo della festa il vino viene meno,
non hanno più la gioia, la danza finirà.
La Madre dice ai servi: «Udite la Parola
che il Figlio mio vi dona, lui vi disseterà».
\endverse
\beginverse
Si arrossano le coppe di vino nuovo colme,
è il dono dello Sposo per nuova fedeltà.
Danziamo allor fratelli: del Re noi siamo figli,
a lui cantiamo lodi per sempre, alleluja!
\endverse
\beginverse*
Danzano con gioia per il figlio del re. \rep{3}
\endverse
\endsong

% Fonte: LibrettoAccordi.pdf (Canti Coro), K8.
% Nell'elenco: Finale, Ordinario (due righe). CD: Cerco il tuo volto (2000).
\beginsong{Resta accanto a me}[by={Gen Verde (CD Cerco il tuo volto)}]
\beginchorus
Ora vado sulla mia strada, con l'amore tuo che mi guida.
O Signore, ovunque io vada resta accanto a me!
Io ti prego, stammi vicino ogni passo del mio cammino.
Ogni notte, ogni mattino, resta accanto a me.
\endchorus
\beginverse
Il tuo sguardo puro sia luce per me
e la tua parola sia voce per me.
Che io trovi il senso del mio andare solo in te,
nel tuo fedele amare il mio perché.
\endverse
\beginverse
Fa' che chi mi guarda non veda che te,
fa' che chi mi ascolta non senta che te.
E chi pensa a me, fa' che nel cuore pensi a te
e trovi quell'amore che hai dato a me.
\endverse
\endsong

% Fonte: LibrettoAccordi.pdf (Canti Coro), K2.
% Nell'elenco: Finale, Ordinario. CD: Una storia che cambia (1984).
\beginsong{Resta qui con noi}[by={Gen Rosso (CD Una storia che cambia)}]
\beginverse
Le ombre si distendono, scende ormai la sera
e s'allontanano dietro i monti
i riflessi di un giorno che non finirà,
di un giorno che ora correrà sempre,
perché sappiamo che una nuova vita
da qui è partita e mai più si fermerà.
\endverse
\beginchorus
Resta qui con noi, il sole scende già,
resta qui con noi, Signore, è sera ormai.
Resta qui con noi, il sole scende già,
se tu sei fra noi la notte non verrà.
\endchorus
\beginverse
S'allarga verso il mare il tuo cerchio d'onda,
che il vento spingerà fino a quando
giungerà ai confini di ogni cuore,
alle porte dell'amore vero.
Come una fiamma che dove passa brucia,
così il tuo amore tutto il mondo invaderà.
\endverse
\beginverse
Davanti a noi l'umanità lotta, soffre e spera
come una terra che nell'arsura
chiede l'acqua da un cielo senza nuvole,
ma che sempre le può dare vita.
Con te saremo sorgente di acqua pura,
con te fra noi il deserto fiorirà.
\endverse
\endsong

% Fonte: LibrettoAccordi.pdf (Canti Coro), K4 (li' senza autore).
% Nell'elenco: Finale (autore Baggio).
\beginsong{Siamo Chiesa del Signore}[by={F. Baggio}]
\beginchorus
Siamo Chiesa del Signore, siamo l'unica famiglia di Dio!
Uno spirito, un solo corpo, siamo tutti comunità.
\endchorus
\beginverse
Nella casa del Padre soli non siamo mai: cittadini del Regno, familiari di Dio,
con gli apostoli edificati sullo stesso Gesù.
\endverse
\beginverse
E la pietra angolare è il Signore Gesù: costruzione ordinata, tempio santo di Dio,
dolce dimora dello Spirito sino all'eternità.
\endverse
\beginverse
Una sola speranza nel Signore Gesù: divenire conformi all'immagine sua,
è chiamata alla santità tutta l'umanità.
\endverse
\endsong

% Fonte: "10 - Finali/vieni e seguimi Sol.pdf" e "Vieni e Seguimi (La).pdf".
% Nell'elenco: senza tipologia/periodo (il file e' tra i Finali).
\beginsong{Vieni e seguimi}[by={Gen Rosso}]
\beginverse
Lascia che il mondo vada per la sua strada.
Lascia che l'uomo ritorni alla sua casa.
Lascia che la gente accumuli la sua fortuna.
Ma tu, tu vieni e seguimi, tu, vieni e seguimi.
\endverse
\beginverse
Lascia che la barca in mare spieghi la vela.
Lascia che trovi affetto chi segue il cuore.
Lascia che dall'albero cadano i frutti maturi.
Ma tu, tu vieni e seguimi, tu, vieni e seguimi.
\endverse
\beginchorus
E sarai luce per gli uomini
e sarai sale della terra
e nel mondo deserto aprirai
una strada nuova. \rep{2}
\endchorus
\beginverse*
E per questa strada va', va',
e non voltarti indietro, va',
e non voltarti indietro.
\endverse
\endsong

% Fonte: LibrettoAccordi.pdf (Canti Coro), A7 (li' tra i canti d'ingresso).
% Nell'elenco: Finale.
\beginsong{Vivere la vita}[by={Gen Verde}]
\beginverse
Vivere la vita con le gioie e i dolori di ogni giorno,
è quello che Dio vuole da te.
Vivere la vita e inabissarti nell'amore è il tuo destino,
è quello che Dio vuole da te.
Fare insieme agli altri la tua strada verso lui,
correre con i fratelli tuoi.
Scoprirai allora il cielo dentro di te,
una scia di luce lascerai.
\endverse
\beginverse
Vivere la vita è l'avventura più stupenda dell'amore,
è quello che Dio vuole da te.
Vivere la vita è generare ogni momento il paradiso:
è quello che Dio vuole da te.
Vivere perché ritorni al mondo l'unità,
perché Dio sta nei fratelli tuoi.
Scoprirai allora il cielo dentro di te,
una scia di luce lascerai. \rep{2}
\endverse
\endsong

% Fonte: LibrettoAccordi.pdf (Canti Coro), J24 (li' tra i canti di comunione).
% Nell'elenco: Finale, Ordinario. CD: In cerca d'autore (1977).
\beginsong{Vocazione}[by={P. A. Sequeri (CD In cerca d'autore)}]
\beginverse
Era un giorno come tanti altri, e quel giorno lui passò.
Era un giorno come tutti gli altri e, passando, mi chiamò.
Come lo sapesse che il mio nome era proprio quello,
come mai vedesse proprio me nella sua vita non lo so.
Era un giorno come tanti altri, e quel giorno mi chiamò.
\endverse
\beginchorus
Tu, Dio, che conosci il nome mio, fa' che ascoltando la tua voce
io ricordi dove porta la mia strada, nella vita all'incontro con te. \rep{2}
\endchorus
\beginverse
Era l'alba triste e senza vita, e qualcuno mi chiamò:
era un uomo come tanti altri, ma la voce, quella no.
Quante volte un uomo con il nome giusto mi ha chiamato,
una volta sola l'ho sentito pronunciare con amore.
Era un uomo come nessun altro, e quel giorno mi chiamò.
\endverse
\endsong

\end{songs}
